\section{Reliability: will it be the same tomorrow, at 3 a.m.?}
\label{sec:reliability}

The clinician needs the same answer tomorrow as today. The sixth question asks
whether the system behaves consistently over time and load, how model drift is
detected and handled, and how uncertainty is quantified and communicated.
Reliability is not the same as accuracy: a system can be accurate on the day it is
validated and quietly degrade as the population, the instruments, or the model
itself shift. The engineering discipline for this is a formal credibility
assessment, and the framework follows the structure of ASME V\&V 40, which grades
how much trust a computational model has earned for a stated use and ties that
trust to the consequences of being wrong \cite{asme2018vv40}.

The mechanism turns that discipline into running controls. Uncertainty
quantification, the third leg of the ten-gate VVUQ examination, attaches a
confidence bound to every estimate, and the gate compares that bound against a
pre-set tolerance before the action proceeds: an estimate whose uncertainty is too
wide does not pass quietly, it takes the ESCALATE path. Around that point-of-
decision check sit three further controls: a reproducibility check that confirms
the same input still yields the same output, a drift monitor that watches inputs
and performance for a statistically significant shift, and a scheduled
revalidation that re-earns credibility on a fixed cadence. Together these implement
the adoption action of continuous monitoring, so reliability is maintained rather
than assumed after go-live.

\autoref{tab:reliability} sets out each control, the trigger that fires it, and the
response it produces. Figure 8 shows the point-of-decision case: an estimate
carries a confidence interval to the gate, which asks whether the estimate sits
within bound and resolves to report or escalate accordingly. The author's
verification-automation lineage demonstrates this at scale, where a code-generation
process was held to account against the FDA FAERS adverse-event record rather than
against its own internal claims \cite{Kawchak2025CodeGenerationCompetition16},
which is the external check that keeps reproducibility honest.

The phrase ``at 3 a.m.'' is the real test, because reliability that holds only
during business hours is not reliability. The night-shift case must be treated
exactly like noon: the same gates, the same uncertainty bounds, the same drift
monitor, and the same escalation to a qualified human if a bound is exceeded. The
illustrative H. R. 9510 would make scheduled revalidation and drift monitoring
standing obligations of the trial rather than optional maintenance
\cite{Kawchak2026HR9510Bill}, so that a clinician arriving at 3 a.m. inherits the
same earned credibility as the team that signed off at noon.

\begin{cfwfig}
\node[cfone] (a) at (0,0) {Estimate}; \node[cftwo] (b) at (3.2,0) {Confidence interval}; \node[cfdec] (c) at (6.8,0) {Within bound?};
\node[cfhope] (d) at (10.2,0.8) {Report}; \node[cfthree] (e) at (10.2,-0.8) {Escalate};
\draw[cfedge] (a)--(b); \draw[cfedge] (b)--(c); \draw[cfedge] (c)--(d); \draw[cfedge] (c)--(e);
\end{cfwfig}
\vspace{-0.2cm}
\cfwcaption{Figure 8. The uncertainty gate (Mermaid 15 reproduced as colored TikZ):
an estimate carries a confidence interval to the gate, which reports the result if
it is within bound and escalates to a qualified human if it is not.}

\needspace{9\baselineskip}\begin{table}[H]
\footnotesize
\begin{tabularx}{\textwidth}{@{}L{3.2cm} Y L{3.2cm}@{}}
\toprule
\textbf{Control} & \textbf{Trigger} & \textbf{Response}\\
\midrule
Reproducibility check & The same input yields a different output than recorded & Halt and investigate before further use\\
Drift monitor & Input distribution or performance shifts beyond a set threshold & Flag the drift and open continuous monitoring\\
Uncertainty bound & An estimate's confidence interval exceeds the gate tolerance & Escalate the case to a qualified human\\
Scheduled revalidation & The fixed cadence or a triggering drift event is reached & Re-run the credibility assessment and gates\\
\bottomrule
\end{tabularx}
\caption{Reliability controls: the trigger that fires each control and the response
it produces, implementing continuous monitoring after go-live.}
\label{tab:reliability}
\end{table}

Reliability is answered by honesty about doubt: a quantified confidence interval, a
plan for drift, and a guarantee that 3 a.m. is treated exactly like noon. The
uncertainty gate, the drift monitor, and scheduled revalidation are the concrete
controls that let a clinician trust the system will give the same defensible answer
tomorrow, on the night shift, as it gave the day it was validated.
